% !TEX root = ../main.tex
\section{Symmetry data and checks}\label{app:symmetry-data-and-checks}

This appendix gives the finite symmetry data and calculations used in
\cref{sec:symmetry-reduced-domain}. The symmetry actions, representative
selection and resulting domain inequalities are stated and explained there.
All finite comparisons below use exact arithmetic in $\mathbb Q(\sqrt5)$.

\subsection{Quaternion families and the full symmetry group}\label{app:quaternion-families-and-the-full-symmetry-group}

The set $\mathcal Q$ defined in
\cref{subsec:symmetries-preserve-the-projection-problem} consists of the
unit quaternions in these three families:
\[
\begin{array}{lr}
  \text{all coordinate permutations of }(\pm1,0,0,0), & 8,\\[2pt]
  \frac12(\pm1,\pm1,\pm1,\pm1), & 16,\\[2pt]
  \text{even coordinate permutations of }
    \frac12(0,\pm1,\pm\varphi,\pm(-1+\varphi)), & 96.
\end{array}
\]
Signs are chosen independently. There are 120 quaternions altogether: the
unit icosians, which form the binary icosahedral group
\citep{ConwaySmith2003}.
The two quaternions in each pair $g,-g$ represent the same rotation.
Composing this rotation with central inversion gives an orientation-reversing
symmetry. The 60 quaternion pairs therefore give 60 rotations in $\mathcal G$
and 60 orientation-reversing symmetries in $-\mathcal G$.

The finite calculations behind this description use the exact vertex
coordinates from \cref{eq:rhombicosidodecahedron-vertices}. They verify
that every listed quaternion has norm one, that the represented rotations
are distinct up to quaternion sign and closed under composition, and
that each takes every vertex to another vertex. All these checks use exact
arithmetic in $\mathbb Q(\sqrt5)$.

\paragraph{Completeness of the symmetry list.}
To see that there are no further body symmetries, fix an ordered triangle
formed by three vertices at pairwise distance two. Exact enumeration gives
20 unordered triangles, with each vertex belonging to exactly one, and verifies
that the three position vectors of each triangle are linearly independent.
A body symmetry must take the chosen triangle to one of them. There are
60 choices for its first image vertex and two for its second; the third is
then fixed. The three independent images determine the linear map, so
there are at most $60\cdot2=120$ orthogonal body symmetries. The 60
rotations and their products with central inversion attain this bound.
Of the 60 orientation-reversing symmetries, 15 are reflections; the whole
group is the full icosahedral group, isomorphic to $A_5\times\{\pm1\}$.

\subsection{Reflection normals and the viewing cone}\label{app:reflection-normals-and-the-viewing-cone}

The fifteen normal lines come from the quaternions in $\mathcal Q$ whose
scalar part is zero. Orient their unit normals $\mathbf n$ toward
$\mathbf h=(1,2,9)$ as in \cref{subsec:rotation-and-reflection-reduction}.
The intersection of their closed halfspaces $\mathbf n\cdot\mathbf u\ge0$ is the cone in
\cref{eq:viewing-cone}. This identification needs only two finite checks
on the listed normals.
For the inclusion of the intersection in the cone, note that the three
planes bounding the cone are among the fifteen, with inward normals
proportional to $(1,0,0)$, $(0,1,0)$ and $(-\varphi,-\varphi^2,1)$.
For the reverse inclusion, every oriented normal has nonnegative dot product with each of
the three vectors
\[
  (0,0,1),\qquad(1,0,\varphi),\qquad(0,1,\varphi^2).
\]
These vectors generate the whole cone: a vector satisfying
\cref{eq:viewing-cone} can be written as
\[
  \mathbf u=u_1(1,0,\varphi)+u_2(0,1,\varphi^2)
    +(u_3-\varphi u_1-\varphi^2u_2)(0,0,1),
\]
with all coefficients nonnegative, so every vector of the cone lies in all
fifteen halfspaces.
The normal orientations and these dot products are checked exactly in
$\mathbb Q(\sqrt5)$.

\subsection{Selected quaternions for the rotation comparisons}\label{app:selected-quaternions-for-the-rotation-comparisons}

The representative $(1,\mathbf r)$ has the largest scalar part among its
right multiples by elements of $\mathcal Q$, which gives the inequality
\[
  \operatorname{scalar}((1,\mathbf r)g)\le1
\]
for every $g\in\mathcal Q$. The domain $D$ uses only twelve of them, those
in \cref{eq:rotation-comparisons}; the others hold as well, but $D$ needs
only to contain a representative of every configuration.
Use cyclic subscripts, with $r_4=r_1$.
For each cyclic pair $i,i+1$ and signs $\epsilon,\delta\in\{-1,1\}$,
$\mathcal Q$ contains a quaternion with scalar part $\varphi/2$ and vector
coordinates $-\epsilon/2$ and $-(-1+\varphi)\delta/2$ in those two
positions, with the remaining coordinate zero. Substituting these twelve
quaternions gives precisely \cref{eq:rotation-comparisons}.
